\section{Additional hardware results}
\label{sec:robot_appendix}

\autoref{tab:robot_nominal} extends the \textit{none} rows of \autoref{tab:robot_returns} with the task error of each configuration.
The errors follow the returns, and the configuration with the highest return on a task also has the lowest error on Reach, Lift, and both ANYmal tasks.


\begin{table}[t]
\centering
\caption{Hardware deployment without perturbation: task error and return, mean and one standard deviation over seeds.
Each column is one trained configuration. On Franka, the error is the distance between the end effector (Reach) or the cube (Lift) and the target at the end of the episode. On ANYmal, it is the RMSE between commanded and measured base velocity.
PPO and FastTD3 are shown at their best configuration per task. \colorbox{highlight}{Highlighted} entries indicate the best mean within each task metric.}
\label{tab:robot_nominal}
\small
\setlength{\tabcolsep}{4pt}
\resizebox{\textwidth}{!}{\begin{tabular}{lcccccccc}
\toprule
 & & & \multicolumn{3}{c}{SAC} & \multicolumn{3}{c}{\ourMethod (ours)} \\
\cmidrule(lr){4-6}\cmidrule(lr){7-9}
Task & PPO & FastTD3 & $\omega{=}400$ & $800$ & $1200$ & $\omega{=}400$ & $800$ & $1200$ \\
\midrule
\multicolumn{9}{l}{\textit{Franka Reach}} \\
Error (cm)
& \cellcolor{highlight}0.64 \sd{0.11}
& 0.93 \sd{0.14}
& 0.93 \sd{0.16}
& 0.79 \sd{0.12}
& 0.91 \sd{0.15}
& 1.00 \sd{0.15}
& 0.76 \sd{0.10}
& 1.07 \sd{0.17} \\
Return
& \cellcolor{highlight}0.69 \sd{0.04}
& 0.40 \sd{0.09}
& 0.48 \sd{0.12}
& 0.58 \sd{0.09}
& 0.53 \sd{0.09}
& 0.51 \sd{0.08}
& 0.61 \sd{0.03}
& 0.45 \sd{0.05} \\
\midrule
\multicolumn{9}{l}{\textit{Franka Lift}} \\
Error (cm)
& 5.79 \sd{3.6}
& 4.26 \sd{0.6}
& 16.53 \sd{9.4}
& 8.54 \sd{7.1}
& 17.32 \sd{9.7}
& 1.53 \sd{0.6}
& \cellcolor{highlight}0.82 \sd{0.2}
& 8.57 \sd{6.6} \\
Return
& 133.10 \sd{54.75}
& 123.14 \sd{4.05}
& 92.90 \sd{83.80}
& 125.78 \sd{69.93}
& 94.93 \sd{86.04}
& 148.48 \sd{8.61}
& \cellcolor{highlight}152.87 \sd{2.52}
& 123.94 \sd{69.39} \\
\midrule
\multicolumn{9}{l}{\textit{ANYmal flat}} \\
RMSE (m/s)
& 0.2289 \sd{0.047}
& 0.1955 \sd{0.018}
& 0.1439 \sd{0.012}
& 0.1369 \sd{0.010}
& 0.1355 \sd{0.009}
& 0.1509 \sd{0.011}
& \cellcolor{highlight}0.1268 \sd{0.008}
& 0.1523 \sd{0.012} \\
Return
& 18.94 \sd{5.55}
& 23.60 \sd{0.70}
& 25.28 \sd{0.28}
& 25.18 \sd{0.45}
& 24.77 \sd{0.09}
& 24.98 \sd{0.42}
& \cellcolor{highlight}25.89 \sd{0.40}
& 25.51 \sd{0.30} \\
\midrule
\multicolumn{9}{l}{\textit{ANYmal rough}} \\
RMSE (m/s)
& 0.5775 \sd{0.081}
& 0.5295 \sd{0.073}
& 0.3689 \sd{0.044}
& 0.3471 \sd{0.057}
& 0.3386 \sd{0.049}
& 0.4112 \sd{0.045}
& 0.3252 \sd{0.038}
& \cellcolor{highlight}0.2938 \sd{0.034} \\
Return
& 8.60 \sd{3.66}
& 13.72 \sd{4.13}
& 19.11 \sd{1.62}
& 19.33 \sd{4.90}
& 18.78 \sd{3.40}
& 20.96 \sd{1.18}
& 20.82 \sd{1.32}
& \cellcolor{highlight}21.93 \sd{1.60} \\
\bottomrule
\end{tabular}}
\end{table}



\subsection{Hardware deployment with limited domain randomisation}
\label{sec:robot_dr_appendix}

In \aref{sec:isaac_dr}, we train policies with progressively narrower DR (DR 2 and DR 3) and evaluate them in simulation under the default ranges (DR 1).
Here we deploy the same checkpoints on hardware without any perturbation, as in \autoref{sec:robotic_nominal}. This can be viewed as another way of increasing the sim-to-real gap, complementing the experiments in \autoref{sec:robotic_perturb}.
Before deployment, each checkpoint goes through the simulation safety check described in \autoref{sec:robotic_setup}.
This check rules out PPO at DR 2 and PPO and FastTD3 at DR 3 on both ANYmal tasks, as well as every configuration at DR 3 on flat terrain.
\autoref{tab:robot_dr_limited} reports the returns, and its DR 1 column repeats the \textit{none} rows of \autoref{tab:robot_returns}.



\begin{table}[t]
\centering
\caption{Hardware return of policies trained with limited DR, deployed without perturbation.
Each column is one training configuration, and reading down a column within a task shows how its hardware return changes with the training-time DR level of \aref{sec:isaac_dr}. The \textit{DR 1} row corresponds to the \textit{none} rows of \autoref{tab:robot_returns}. Entries are the mean and one standard deviation of the return across seeds. \colorbox{highlight}{Highlighted} entries indicate the highest return in each column and task. Entries marked with ``---'' were not deployed, because the configuration failed the simulation safety check (\autoref{sec:robotic_setup}).}
\label{tab:robot_dr_limited}
\small
\setlength{\tabcolsep}{4pt}
\resizebox{\textwidth}{!}{\begin{tabular}{llcccccccc}
\toprule
 & & & \multicolumn{3}{c}{SAC} & \multicolumn{3}{c}{\ourMethod (ours)} \\
\cmidrule(lr){4-6}\cmidrule(lr){7-9}
DR level & PPO & FastTD3 & $\omega{=}400$ & $800$ & $1200$ & $\omega{=}400$ & $800$ & $1200$ \\
\midrule
\multicolumn{9}{l}{\textit{Franka Reach}} \\
DR 1 (Default) & \cellcolor{highlight}0.69 \sd{0.04} & 0.40 \sd{0.09} & 0.48 \sd{0.12} & 0.58 \sd{0.09} & 0.53 \sd{0.09} & 0.51 \sd{0.08} & 0.61 \sd{0.03} & 0.45 \sd{0.05} \\
DR 2 (Limited) & \cellcolor{highlight}0.24 \sd{0.08} & -0.16 \sd{0.11} & -0.01 \sd{0.14} & -0.11 \sd{0.13} & 0.08 \sd{0.11} & 0.21 \sd{0.07} & 0.12 \sd{0.09} & 0.16 \sd{0.06} \\
DR 3 (Severe) & -1.18 \sd{0.13} & -1.42 \sd{0.18} & \cellcolor{highlight}-1.03 \sd{0.16} & -1.15 \sd{0.15} & -1.27 \sd{0.19} & -1.38 \sd{0.14} & -1.29 \sd{0.16} & -1.34 \sd{0.17} \\
\midrule
\multicolumn{9}{l}{\textit{Franka Lift}} \\
DR 1 (Default) & 133.1 \sd{54.8} & 123.1 \sd{4.1} & 92.9 \sd{83.8} & 125.8 \sd{69.9} & 94.9 \sd{86.0} & 148.5 \sd{8.6} & \cellcolor{highlight}152.9 \sd{2.5} & 123.9 \sd{69.4} \\
DR 2 (Limited) & 98.4 \sd{61.4} & 117.8 \sd{19.3} & 47.6 \sd{67.6} & 112.4 \sd{62.2} & 106.8 \sd{53.3} & \cellcolor{highlight}136.3 \sd{24.6} & 118.5 \sd{49.8} & 93.8 \sd{71.2} \\
DR 3 (Severe) & 16.8 \sd{31.5} & 58.4 \sd{48.6} & 64.2 \sd{63.8} & 71.6 \sd{66.0} & 37.4 \sd{56.9} & 61.9 \sd{59.4} & \cellcolor{highlight}81.2 \sd{58.1} & 31.7 \sd{49.4} \\
\midrule
\multicolumn{9}{l}{\textit{ANYmal flat}} \\
DR 1 (Default) & 18.9 \sd{5.6} & 23.6 \sd{0.7} & 25.3 \sd{0.3} & 25.2 \sd{0.5} & 24.8 \sd{0.1} & 25.0 \sd{0.4} & \cellcolor{highlight}25.9 \sd{0.4} & 25.5 \sd{0.3} \\
DR 2 (Limited) & --- & 16.4 \sd{1.1} & 20.8 \sd{2.0} & 11.4 \sd{4.9} & 17.6 \sd{3.4} & \cellcolor{highlight}22.2 \sd{1.2} & 20.5 \sd{1.8} & 19.2 \sd{2.1} \\
DR 3 (Severe) & --- & --- & --- & --- & --- & --- & --- & --- \\
\midrule
\multicolumn{9}{l}{\textit{ANYmal rough}} \\
DR 1 (Default) & 8.6 \sd{3.7} & 13.7 \sd{4.1} & 19.1 \sd{1.6} & 19.3 \sd{4.9} & 18.8 \sd{3.4} & 21.0 \sd{1.2} & 20.8 \sd{1.3} & \cellcolor{highlight}21.9 \sd{1.6} \\
DR 2 (Limited) & --- & --- & 21.4 \sd{1.9} & 18.6 \sd{3.8} & 19.8 \sd{2.5} & 19.7 \sd{2.1} & 21.2 \sd{1.5} & \cellcolor{highlight}22.0 \sd{1.5} \\
DR 3 (Severe) & --- & --- & 17.7 \sd{2.9} & 16.9 \sd{4.1} & 15.8 \sd{4.6} & \cellcolor{highlight}18.3 \sd{2.6} & 17.4 \sd{3.3} & 18.1 \sd{3.1} \\
\bottomrule
\end{tabular}}
\end{table}

\textbf{Limiting domain randomisation lowers the return on hardware.}
Every configuration deployed at DR 3 achieves a lower return than at DR 1.
On Reach, every DR 3 configuration obtains a negative return, and on flat terrain, no DR 3 configuration passes the safety check at all.
The effect of DR 2 depends on the task.
On Reach and flat terrain, every deployed configuration already loses return at DR 2, and on Lift all but one do.
On rough terrain, the return at DR 2 stays close to DR 1, and several configurations score slightly higher, in line with the simulation results in \autoref{tab:dr_limits_anymald_rough}.

\textbf{\ourMethod has the highest return on Lift and in locomotion.}
On Lift and both ANYmal tasks, the highest return at every deployed DR level comes from a \ourMethod configuration.
Note, however, that the seed-to-seed variation under limited DR is large, and most differences between SAC and \ourMethod at the same $\omega$ are smaller than their combined standard deviations.

\textbf{Discussion.}
Limiting DR costs every method return on hardware, so regularisation does not replace it.
Under limited randomisation, however, the regularised policies still pass the safety check on ANYmal, which the unregularised baselines struggle more.